\section{Related Works}
\label{sec:related_works}

Most annotated resources for affective computing focus on English. While multilingual initiatives like BRIGHTER \cite{brighter_project} aggregate emotion data across 28 languages and 100,000 instances, they omit Italian, pointing users toward standalone external sets like FEEL-IT \cite{bianchi2021feel}. This gap leaves Italian social media underrepresented in multi-label benchmarks, despite distinct cultural and linguistic norms in regional political discourse.

Early Italian resources focused strictly on binary or three-class polarity. The first dedicated shift toward fine-grained categories came with \textit{FEEL-IT} \cite{bianchi2021feel}, which collected 2,037 tweets across politics, sports, and healthcare during 2020. Annotated with Ekman's four basic emotions (anger, joy, sadness, fear), FEEL-IT established that Italian language models like UmBERTo \cite{musixmatch-2020-umberto} produce reliable affective classifiers. Yet a four-class schema misses mixed feelings and conversational nuances common in viewer reactions. To address broader emotional ranges, \textit{MultiEmotions-It} \cite{MultiEmotion} annotated YouTube and Facebook comments across Plutchik's eight primary emotions \cite{Plutchik}, polarity, and sarcasm. Their findings showed that while polarity models perform consistently, multi-label emotion recognition and sarcasm detection suffer from severe class sparsity and label ambiguity.

Second-screen platforms have changed how broadcasters analyze viewer engagement \cite{SocialTV}. Traditional television ratings (like Auditel) measure raw viewer counts; they cannot gauge whether an audience is laughing, outraged, or sympathetic. The \textit{M.EMO.RAI} project was formed to address this limitation, establishing a joint partnership between the University of Turin and Rai (Radiotelevisione Italiana) \cite{memorai}. The goal was an automated system to track audience reactions during live broadcasts by ingesting public social network commentary.

This partnership led to the \textit{EMiT} (Emotions in Italian) corpus, benchmarked as an EVALITA 2023 shared task \cite{emit}. With 5,966 social media posts from live broadcasts (including the Sanremo Festival and scripted drama), EMiT adopted Plutchik's taxonomy. It also introduced the \textit{Target} dimension, distinguishing whether a comment attacks the broadcast \textit{Content} (the topic, guest, or political issue) or evaluates the \textit{Container} (the host, studio format, or network production). EVALITA benchmark results showed that while classifiers perform well on familiar shows, accuracy drops on out-of-domain broadcasts. We use EMiT as an external textual corpus to pre-train and enrich our baseline emotion encoders.

Irony detection requires models to recognize when a speaker's intended meaning opposes their literal words. In Italian NLP, the standard benchmark is \textit{IronITA} \cite{ironita2018}, curated for EVALITA 2018 with nearly 5,000 tweets annotated for irony and sarcasm (defined as hostile or denigratory mockery). Results on IronITA showed that text-only classifiers struggle on short posts when pragmatic context is missing. To provide our models with domain-adapted pragmatic competence, we use IronITA to fine-tune UmBERTo before evaluating on television commentary.

Linguistic text only tells half the story in modern social media. In their study on document intent, Kruk et al. \cite{MultiM} demonstrated that images in social posts rarely just illustrate captions; they alter communicative intent through distinct semiotic relationships, ranging from complementary addition to direct divergence. Incongruent or divergent pairings generate satirical humor and memes, where the image subverts the text. Kruk et al. studied English Instagram posts. The semiotic relationship between text and visual media in Italian political satire and reality television remains largely unstudied. MultiEMiT applies this semiotic framework directly to Italian Social TV.
